\textbf{Problem 4.1 (Dujella, 2026).} Let $k \geqslant 2$ be an integer. Is it true that the Pellian equation
\[
x^2 - (k^2+1)y^2 = k^2
\]
has at most one solution with $0<y<k-1$ and $x>0$?

\medskip

\textit{Remark.} Matthews, Robertson and White (Glas. Mat. Ser. III \textbf{48} (2013), 265--289) verified that the statement is true for $k \leqslant 2^{50}$, while Le and Srinivasan (Glas. Mat. Ser. III \textbf{58} (2023), 59--65) proved it for all $k$ of the form $k=p^mq^n$, where $p,q$ are distinct odd primes and $m,n$ are positive integers.
